\section{Interpretación de Benchmark y direcciones de investigación}

\subsection{Interpretación de las curvas de Arc AGI / O3}
Arc AGI, como indicador central de un reasoning benchmark, permitió inicialmente que un presupuesto de 1,000 dólares alcanzara 87.5\% de accuracy, mientras que por debajo de 20 dólares solo se mantenía en 25\%. Chen señala que esta curva no es simplemente cuestión de «pagar para lograrlo», sino que requiere coordinar al mismo tiempo sampler, search y verifier, para que cada partida de presupuesto se convierta en accuracy efectiva. La práctica de demos como OpenAI O3, AlphaProof/AlphaGeometry, depende firmemente del inference-time compute, y no del dominio de entrenamiento.

\begin{table}[H]
\centering
\caption{Comparación entre presupuesto y desempeño en tiempo de inferencia de los principales Benchmark}
\begin{tabular}{p{0.30\textwidth}p{0.28\textwidth}p{0.32\textwidth}}
\toprule
Benchmark & Token / Cost budget & Observed behavior \\
\midrule
Arc AGI & \$20~\$1,000 per query & 25\% → 87.5\% accuracy; the curve saturates after \$500 without better search/verifier \\
OpenAI O3 & \$20~\$200 & Consistency + trace feedback gives human-level accuracy once budget crosses \$100 \\
AlphaProof / AlphaGeometry & 50~200 token steps & RL-based reward tweaks inference-time policy to maintain correctness under limited latency \\
\bottomrule
\end{tabular}
\end{table}

\subsection{Tablero de Benchmark y niveles de presupuesto}
Chen recomienda construir para cada benchmark un tablero de ``budget bin'': dividir el presupuesto de token/cost en tres niveles, Conservative, Balanced y Aggressive, que corresponden respectivamente a los casos de `20~80 token`, `80~300 token` y `>300 token`. Cada bin debe registrar accuracy, stepwise verifier margin y token per vote, para poder ubicar rápidamente en tiempo de ejecución «en qué bin se encuentra el estado actual». Esta estratificación permite que el equipo, al publicar los resultados del benchmark hacia afuera, diga directamente «en el bin Balanced, nuestra accuracy es de 87.5\% pero la latency es de 3.4s».

\begin{table}[H]
\centering
\caption{Ejemplo de tablero de Benchmark}
\begin{tabular}{p{0.26\textwidth}p{0.30\textwidth}p{0.30\textwidth}}
\toprule
Bin & Configuración típica & Indicadores de monitoreo \\
\midrule
Conservative & $\leq 80$ token, $\leq 2$ iterations & Accuracy, sampling diversity, latency \\
Balanced & 80~300 token, 2~3 iterations & Verifier margin, branching ratio, token per candidate \\
Aggressive & > 300 token, 3+ iterations & Score per token, trace iterations, SLO breach rate \\
\bottomrule
\end{tabular}
\end{table}

\begin{importantbox}{Recomendaciones para implementar el tablero de Benchmark}
1. Configurar para cada bin un `alert level`, y cambiar de bin cuando la latency supere el target o el verifier margin caiga por debajo del threshold; 2. hacer push del state del bin al dashboard, para que el equipo de operaciones pueda leer rápidamente «en qué nivel está corriendo el benchmark hoy»; 3. combinar esto con el monitoreo de score-per-token dentro del bin, para detectar a tiempo el early warning de que el token budget se está saliendo de control.
\end{importantbox}

\subsection{Preguntas de investigación abiertas}
\begin{itemize}
\item ¿Cómo ampliar el search depth y el iteration depth con compute limitado en escenarios de baja latency?
\item ¿Es posible diseñar un «budget-aware verifier» que reporte en tiempo real el tradeoff entre token y candidate, sin esperar a que termine la inference?
\item ¿Cómo garantizar la generalidad del trace feedback en el reasoning multimodal o de imagen más texto?
\item ¿Cómo, en un pipeline de despliegue real, encadenar indicadores como sampling diversity, stepwise scaling metric y trace drift en un SLO ejecutable?
\end{itemize}

\begin{itemize}
\item ¿Puede un pipeline de inferencia ajustar de manera adaptativa el branching ratio y el iteration depth mediante learning-to-rank, en lugar de fijar umbrales manualmente?
\item Para un adversarial prompt, ¿cómo combinar el verifier margin con el prompt-editing loop, para evitar que un razonamiento erróneo se refuerce mediante el feedback?
\item En el reasoning multimodal, ¿qué incertidumbres persisten en la correspondencia entre el presupuesto de token y el vision token length?
\end{itemize}

Estas preguntas se mencionaron repetidamente en el Q\&A posterior a las 01:19:00, y Chen sugiere explorarlas mediante el enfoque «metric-first hypothesis»: primero proponer una metric central (como `score per token`) y luego iterar el ajuste del pipeline, en lugar de ampliar el token budget guiándose por la intuición.

\begin{knowledgebox}{Asociaciones de direcciones de investigación}
Las preguntas abiertas pueden abordarse desde tres ángulos: budgeting instrumentation, verifier robustness y la ciencia del trace. Chen recomienda escribir los inference-time logs (sampling count/delay/score) en el dashboard, para proporcionar los datos base que permitan en el futuro ajustar automáticamente la estrategia.
\end{knowledgebox}

\subsection{Resumen de la sección}
A través de la interpretación de las curvas de budget/accuracy de Arc AGI, O3 y la serie de DeepMind, vemos que el inference-time compute es a la vez un parámetro de ingeniería y un detonante de direcciones de investigación. En el futuro será necesario mantener una high accuracy bajo low-latency, e integrar los indicadores de trace/verifier/gating en el pipeline de despliegue.

